\documentclass[a4paper,10pt]{article}

% Paquetes
\usepackage{fontspec}
\usepackage{titlesec}
\usepackage{titling}
\usepackage{setspace}
\usepackage[hidelinks]{hyperref}
\usepackage{graphicx}
\usepackage{caption}
\usepackage{subcaption}
\usepackage{float}
\usepackage{xcolor}
\usepackage[affil-it]{authblk}
\usepackage{biblatex}
\addbibresource{Mendeley.bib}

% Idioma
\usepackage[spanish]{babel}
\usepackage[utf8]{inputenc}


% Configuraciones
\setmainfont{Times New Roman}
\hypersetup{
    colorlinks,
    linkcolor={red!50!black},
    citecolor={blue!50!black},
    urlcolor={blue!80!black}
}


\begin{document}
\begin{titlepage}
\centering
\begin{figure}[t]
	\centering
	\includegraphics[scale=0.15]{utn.jpg}
    \vspace{0.5cm}
\end{figure}%
	{\LARGE Proyecto Final\par}
    {\LARGE 2016\par}
	\vspace{1cm}
	{\huge\bfseries Titulo del Proyecto\par}
	\vspace{1cm}
    {\LARGE Tesina\par}
    \vspace{1cm}
	{\Large\itshape Integrante 1, Integrante 2, Integrante 3\par}
	\vfill
\end{titlepage}


\tableofcontents

\newpage
\section{Objeto del Documento}
\label{sec:0_objeto}
Este documento tiene por objeto la explicación  detallada  y  con  el  máximo  rigor  técnico  los conceptos  teóricos  de  aplicación  en  el  proyecto,  con  referencias  bibliográficas  que cubran  los  contenidos  teóricos  clásicos,  la  cual  debe  ser  actualizada  y  de  reconocido prestigio.

\newpage
\section{Abstract}
\label{sec:1_abstract}

\vspace{1cm}
Ejemplo de una imagen cargada:
\begin{figure}[!ht]
  \centering
    \includegraphics[width=0.5\textwidth]{logo.jpg}
      \caption{Este es el logo de la UTN.}
\end{figure}

\vspace{1cm}
Ejemplo de cita:

Todos los cuerpos se mantienen firmes y constantes en su estado de reposo o de movimiento uniforme en línea recta, salvo que se vean forzados a cambiar ese estado por fuerzas impresas.\cite{cita_1}

\vspace{1cm}
Ejemplo de una tabla:
\begin{table}[h]
	\centering
    \begin{tabular}{ | c | c | c | p{5cm} |}
    \hline
    Día & Mes & Año & Comentario \\ \hline
    Lunes & Enero & 01 & Primer dia del año. \\ \hline
    Martes & Enero & 02 & Segundo dia del año. \\ \hline
    \end{tabular}
    \caption{Tabla de los Días}
\end{table}

\newpage
\section{Introducción Teorica del tema a tratar y Marco Teorico}
\label{sec:2_marcoteorico}

\newpage
\section{Desarrollo de la solución a implementar}
\label{sec:3_solucion}

\newpage
\section{Resultados, Mediciones y Verificación}
\label{sec:4_resmedver}

\newpage
\section{Conclusiones}
\label{sec:5_concl}

\newpage
\section{Referencias}
\label{sec:6_refs}

\newpage
\section{Bibliografía}
\label{sec:7_biblio}

\begin{thebibliography}{9}
\bibitem{cita_1}
	NEWTON, Isaac y WHITESIDE, D. T.,
    \textit{"The Mathematical Papers of Issac Newton"}.
    Cambridge University Press,
    UK,
    1971,
    722 p.
    ISBM 978-0-521-07740-8.  
\end{thebibliography}

\newpage
\section{Índice de Figuras y Tablas}
\label{sec:7b_indice}

\listoffigures
\listoftables

\newpage
\section{Apéndices}
\label{sec:8_apendix}

\subsection{Circuitos electricos}
\label{sec:9_circ}

\subsection{Hojas de Datos}
\label{sec:10_hd}

\subsection{Programas Fuente}
\label{sec:11_soft}

\subsection{Layout PCB}
\label{sec:12_pcblay}

\end{document}